% gz-p10-01-1: 切线斜率的几何意义——f(x)=x^2 在 x_0=1 处切线（斜率=2）和割线（斜率→2）
\documentclass[tikz,border=4pt]{standalone}
\usepackage{ctex}
\usepackage{amsmath}
\usepackage{pgfplots}
\pgfplotsset{compat=1.18}

\begin{document}
\begin{tikzpicture}
\begin{axis}[
  width=10cm, height=8cm,
  axis lines=center,
  xlabel={$x$}, ylabel={$y$},
  every axis x label/.style={at={(axis cs:2.6,0)}, anchor=west},
  every axis y label/.style={at={(axis cs:0,6.2)}, anchor=south},
  xmin=-0.5, xmax=2.6,
  ymin=-0.5, ymax=6.2,
  xtick={1,2},
  ytick={1,4},
  tick style={color=black},
  ticklabel style={font=\small},
  clip=true,
]
  % 曲线 y = x^2
  \addplot[black, thick, domain=-0.5:2.5, samples=80] {x^2};
  \node[right, black] at (axis cs:2.0, 4.6) {$y=x^2$};

  % 切点 P(1,1)
  \addplot[only marks, mark=*, mark size=2.2pt, red] coordinates {(1,1)};
  \node[below right, red, font=\small] at (axis cs:1,1) {$P(1,1)$};

  % 切线 y = 2x - 1（斜率 = 2，过 (1,1)）
  \addplot[red, thick, domain=0.2:2.5] {2*x - 1};
  \node[red, font=\small] at (axis cs:2.3, 3.7) {切线$\;k=2$};

  % 割线点 Q(2,4) 与割线 PQ
  \addplot[only marks, mark=*, mark size=2.2pt, blue] coordinates {(2,4)};
  \node[above left, blue, font=\small] at (axis cs:2,4) {$Q(2,4)$};
  \addplot[blue, dashed, thick, domain=0.4:2.4] {3*x - 2};
  \node[blue, font=\small] at (axis cs:0.55, 0.0) {割线$\;k=3$};

  % 注释：Q→P 时割线→切线
  \node[font=\footnotesize, align=left] at (axis cs:0.3, 5.5) {当 $Q\to P$，\\割线斜率 $\to 2$};
\end{axis}
\end{tikzpicture}
\end{document}
